\chapter{Contexto y alcance}

\section{Punto de partida}
El repositorio estaba en el commit \texttt{620aa0c} (``docs: agregar documentacion Tech-Flow''), versión más reciente de \texttt{origin/main} de \texttt{fausto-lang/Tech-Flow}. Se partía de un sistema funcional (inventario, ventas, caja, entradas, garantías) con entorno de pruebas verde y persistencia CSV, sobre el que se pidió una nueva funcionalidad y, a continuación, una ronda de organización del código.

\section{Objetivo}
\begin{enumerate}
\item Añadir un filtro por categorías al catálogo/inventario de las tres pantallas del sistema (administración, almacén y punto de venta), tal como solicitó el docente.
\item Corregir la inconsistencia del precio de venta detectada durante la revisión.
\item Refactorizar puntos duplicados y código muerto detectados al leer el proyecto completo.
\end{enumerate}

\section{Alcance del informe}
Cubre el trabajo realizado entre \texttt{620aa0c} y \texttt{a618180} (``refactor(console): delegate validation to ValidadorDatos''). Los cambios se realizaron en seis commits:

\begin{table}[H]\centering\caption{Commits incluidos en este informe}
\begin{tabular}{lp{.72\textwidth}}\toprule
Commit & Descripción\\\midrule
\texttt{2fe216f} & feat(gui): add category filter to views and product forms\\
\texttt{cabca82} & refactor: share controller UI helpers and centralize CSV file names\\
\texttt{2ecf69d} & refactor(gui): drop Categoria enum, derive filter list from inventory CSV\\
\texttt{7e73fed} & fix(domain): persist sale price using product precioVenta\\
\texttt{1f081df} & refactor(controllers): extract shared save-txt helpers\\
\texttt{a618180} & refactor(console): delegate validation to ValidadorDatos\\\bottomrule
\end{tabular}\end{table}

\section{Estado resultante}
La suite JUnit 5 queda en \textbf{516/516 pruebas en verde} y las cuatro vistas FXML se cargan sin errores (verificación en \chaptername\ 4). El árbol de trabajo quedó limpio y todo el trabajo commiteado.
